\begin{textsample}{POS Dim 1 – vem\_ed\_23 – Score 8.00 – vem\_ed\_23/t124.txt}  \label{ex:f1_pos_141}
Visitante encoraja alunos rumo à Internacionalização Em 15 de março de 2024, a Profa.
Dra.
Cathy Lee Arcuino visitou as Fatecs Ipiranga e São Caetano do Sul.
Na Fatec Ipiranga, ela conversou com uma das \textbf{turmas} do curso pioneiro de Big Data para Negócios, durante a aula do Prof.
Ms.
Carlos Eduardo Dantas de Menezes.
Cathy falou sobre a \textbf{importância} de praticar o inglês, \textbf{expandir} a visão de mundo e interagir com outras culturas.
E \textbf{é} \textbf{possível} fazer isso mesmo sem viajar, por meio dos Intercâmbios Virtuais – nas Fatecs, \textbf{são} conhecidos como Projetos Colaborativos Internacionais (PCIs / Cesu).
Menezes mostrou para a visitante internacional alguns dos resultados dos estudos exploratórios conduzidos por ele com seus alunos de Aprendizado de Máquina sobre as pesquisas de percepção entre alunos e professores envolvidos com os projetos COIL nas Fatecs (\textbf{saiba} mais sobre esses estudos na cobertura da IVEC 2024, nesta edição).
O coordenador do curso de Eventos da Fatec Ipiranga, Marcos Julio, auxiliou a recepção à visitante.
Em seguida, Cathy Lee acompanhou o coordenador dos PCIs / Cesu, Osvaldo Succi Junior, e as professoras Neusa Haruka Sezaki Gritti e Patrícia Sales Patrício (equipe PCIs / Cesu) à Fatec São Caetano, onde \textbf{foram} recebidos pela diretora da unidade, Adriane Monteiro Fontana, \textbf{que} apresentou \textbf{parte} do corpo docente envolvido com PCIs à professora Cathy Lee.

% matched lemmas: expandir, importância, parte, possível, que, saber, ser, turma
\end{textsample}
